\begin{chineseabstract}

随着大语言模型（Large Language Model，LLM）能力提升，
自然语言驱动的Shell命令生成成为降低命令行使用门槛的重要方向。
但通用模型在中文Shell场景仍存在领域知识覆盖不足、命令幻觉与执行风险高等问题。

本文面向中文Shell交互，设计并实现一套检索增强与安全执行一体化智能体系统。
系统采用客户端—服务端架构，以本地量化Qwen-7B为推理核心，
结合基于SQLite的分层记忆机制与Bash/Zsh语法适配，提升多轮交互稳定性。
在检索侧，构建基于ChromaDB与BAAI/bge-small-zh-v1.5的检索增强生成模块，
提出向量相似度与关键词命中率结合的混合重排序方法，
并通过关键词路由将查询分发至命令、安全、任务、模式、示例五类知识库。
在执行侧，提出三级命令安全框架，包含权限分级、
“语法校验—风险评估—效果模拟”预执行流水线与结构化日志；
同时实现面向复杂请求的分步执行机制，支持最多12步任务的逐步生成、校验与确认。

实验结果表明，在80条自然语言查询上，所提检索策略达到Recall@1为0.90、
平均倒数排名（Mean Reciprocal Rank，MRR）为0.9375，
均优于纯向量检索基线，验证了方法有效性与工程可行性。

\chinesekeyword{大语言模型；Shell智能体；检索增强生成；混合检索；命令安全执行}
\end{chineseabstract}
